\documentclass{article}
\usepackage[margin=1in]{geometry}

\usepackage{hyperref}
\usepackage{url}

\usepackage[utf8]{inputenc}
\usepackage[T1]{fontenc}
\usepackage{booktabs}
\usepackage{amsfonts}
\usepackage{nicefrac}
\usepackage{microtype}
\usepackage{xcolor}

\usepackage{amsmath}
\usepackage{amssymb}
\usepackage{mathtools}
\usepackage{amsthm}

\usepackage[inline]{enumitem}
\setitemize{noitemsep,topsep=0pt,parsep=0pt,partopsep=0pt,leftmargin=*}
\setenumerate{noitemsep,topsep=0pt,parsep=0pt,partopsep=0pt,leftmargin=*}

\usepackage[capitalize,noabbrev]{cleveref}
\crefname{section}{\S}{\S\S}

\usepackage{latexsym}
\usepackage{mleftright}
\usepackage{xspace}

\usepackage{stmaryrd}
\usepackage{bbm}

\usepackage{tikz}
\usetikzlibrary{arrows,arrows.meta,shapes,automata,positioning}

%%%%%%%%%%%%%%%%%%%%%%%%%%%%%%%%%%%%%%%%%%%%%%%%%%%%%%%%%%
% MACROS (from paper/macros.tex)
%%%%%%%%%%%%%%%%%%%%%%%%%%%%%%%%%%%%%%%%%%%%%%%%%%%%%%%%%%

\colorlet{MacroColor}{black}
\newcommand{\mymacro}[1]{#1}

\newcommand{\defn}[1]{\textbf{#1}}
\newcommand{\defeq}{\mathrel{\overset{\raisebox{-0.25ex}{\textnormal{\tiny def}}}{=}}}

\newcommand{\card}[1]{\left|#1\right|}
\newcommand{\set}[1]{\left\{ #1 \right\}}

\newcommand{\bigo}{{\mymacro{ \mathcal{O}}}}
\newcommand{\kleene}[1]{#1^*}

%%%%%%%%%%%%%%%%%%%%%%%%%%%%%%%%%%%%%%%%%%%%%%%%%%%%%%%%%%
% PAPER-SPECIFIC NOTATION
%%%%%%%%%%%%%%%%%%%%%%%%%%%%%%%%%%%%%%%%%%%%%%%%%%%%%%%%%%

\definecolor{CharacterColor}{RGB}{98,115,19}
\definecolor{TokenColor}{RGB}{140,10,89}
\definecolor{StormBlue}{RGB}{42,90,120}

\newcommand{\preCover}{{\color{TokenColor}\mathcal{P}}}
\newcommand{\maxQuotient}{{\color{TokenColor}\overline{\mathcal{Q}}}}
\newcommand{\Quotient}{{\color{TokenColor}\mathcal{Q}}}
\newcommand{\Remainder}{{\color{TokenColor}\mathcal{R}}}

\newcommand{\eos}{{\mymacro{\textsc{eos}}}}

\newcommand{\basisSet}[1]{\langle #1\rangle}

\newcommand{\tokenfont}[1]{{\color{TokenColor}\texttt{#1}}}
\newcommand{\tfont}[1]{{\color{TokenColor}\texttt{#1}}}
\newcommand{\charfont}[1]{{\color{CharacterColor}\strut\texttt{#1}}}

\newcommand{\transFnSym}[0]{f}
\newcommand{\transFn}{\ensuremath{{\color{CharacterColor}\transFnSym}}\xspace}

% Source / target symbols and strings
\newcommand{\srcSym}{{\color{TokenColor}x}}
\newcommand{\srcSymp}{{\color{TokenColor}x^\prime}}
\newcommand{\srcStr}{\ensuremath{{\color{TokenColor}\boldsymbol{x}}}\xspace}
\newcommand{\srcStrp}{\ensuremath{{\color{TokenColor}\boldsymbol{x}^\prime}}\xspace}

\newcommand{\tgtSym}{{\color{CharacterColor}y}}
\newcommand{\tgtStr}{{\color{CharacterColor}\boldsymbol{y}}}
\newcommand{\tgtStrp}{{\color{CharacterColor}\boldsymbol{y}^\prime}}

% Alphabets
\newcommand{\srcAlphabet}[0]{{\color{TokenColor}\ensuremath{\mathcal{X}}}\xspace}
\newcommand{\tgtAlphabet}[0]{{\color{CharacterColor}\ensuremath{\mathcal{Y}}}\xspace}
\newcommand{\srcStrings}[0]{\ensuremath{\srcAlphabet^*}\xspace}
\newcommand{\tgtStrings}[0]{\ensuremath{\tgtAlphabet^*}\xspace}

% Source LM
\newcommand{\pSource}{\mymacro{p}}
\newcommand{\prefixSource}{\mymacro{\bar{p}}}
\newcommand{\prefixTarget}{\mymacro{\bar{p}}_{\!\tgtSym}}

% Transducer
\newcommand{\transST}{{\color{CharacterColor}\texttt{f}}}

% States
\newcommand{\States}{\mymacro{\mathsf{S}}}
\newcommand{\InitialStates}[0]{\mymacro{\mathsf{I}}}
\newcommand{\FinalStates}{\mymacro{\mathsf{F}}}

\newcommand{\state}{\mymacro{s}}
\newcommand{\statep}{\mymacro{s^\prime}}

% Cover algorithm
\newcommand{\frontier}{\mathcal{F}}
\newcommand{\powerstate}{S}
\newcommand{\powerstatep}{\powerstate^{\prime}}

\newcommand{\Paths}{\mymacro{\Pi}}
\newcommand{\Path}{\mymacro{\boldsymbol{\pi}}}

\newcommand{\sem}[1]{\llbracket #1 \rrbracket}

\newcommand{\buffer}{{\color{CharacterColor}\boldsymbol{b}}}

% New macros for this report
\newcommand{\committed}{\mathrm{out}}
\newcommand{\lcp}{\mathrm{lcp}}
\newcommand{\dom}{\mathrm{dom}}

%%%%%%%%%%%%%%%%%%%%%%%%%%%%%%%%%%%%%%%%%%%%%%%%%%%%%%%%%%

\theoremstyle{plain}
\newtheorem{theorem}{Theorem}[section]
\newtheorem{proposition}[theorem]{Proposition}
\newtheorem{lemma}[theorem]{Lemma}
\newtheorem{corollary}[theorem]{Corollary}

\theoremstyle{definition}
\newtheorem{definition}[theorem]{Definition}
\newtheorem{example}[theorem]{Example}

\theoremstyle{remark}
\newtheorem{remark}[theorem]{Remark}

\title{Function-Level Characterization\\of Finite Decomposition}
\author{}
\date{}

\begin{document}
\maketitle

\begin{abstract}
The paper's current finite decomposition lemma (Lemma~6.1) gives sufficient
conditions at the \emph{transducer} level: it refers to paths, states,
universality, and finite closure of individual states.  This report develops a
parallel characterization at the \emph{function} level, with no reference to
any particular transducer.  The key concept is the \emph{committed prefix}
$\committed(\srcStr)$---the longest output prefix guaranteed after reading
input $\srcStr$, regardless of future input.  It decomposes the finite
decomposition property into two independent axes: \emph{finite commitment
boundary} (the quotient is finite) and \emph{finite tentative set} (the
remainder is finite), each generalizing strict-prefix monotonicity in a
different direction.
\end{abstract}


\section{Background: The Decomposition is Function-Level}
\label{sec:background}

The quotient and remainder are defined purely in terms of the function
$\transFn \colon \srcStrings \to \tgtStrings$, with no reference to any
transducer.  Recall (paper, Definition~4.1):
%
\begin{align}
  \preCover(\tgtStr) &\defeq \{ \srcStr \in \srcStrings : \tgtStr \preceq \transFn(\srcStr) \}, \\[4pt]
  \maxQuotient(\tgtStr) &\defeq \{ \srcStr \in \srcStrings : \basisSet{\srcStr} \subseteq \preCover(\tgtStr) \},
    \label{eq:maxQ} \\[4pt]
  \Quotient(\tgtStr) &\defeq \mathrm{pr}\!\left(\maxQuotient(\tgtStr)\right),
    \qquad
  \Remainder(\tgtStr) \defeq \preCover(\tgtStr) \setminus \maxQuotient(\tgtStr).
    \label{eq:QR}
\end{align}
%
Here $\basisSet{\srcStr} = \{\srcStrp : \srcStr \preceq \srcStrp\}$ is the
cylinder (upward closure) of $\srcStr$ under the prefix order, and
$\mathrm{pr}(\cdot)$ extracts the minimal prefix-free generating set.

\begin{remark}[Function-level property]
\label{rem:function-level}
The \defn{finite decomposition property}---that $\Quotient(\tgtStr)$ and
$\Remainder(\tgtStr)$ are both finite for every $\tgtStr \in
\tgtStrings$---depends only on $\transFn$, not on how $\transFn$ is
represented.  Two transducers realizing the same function have identical
$\Quotient(\tgtStr)$ and $\Remainder(\tgtStr)$ for every $\tgtStr$.
\end{remark}

This is in contrast to the paper's current Lemma~6.1, whose conditions (safety,
finite closure, universality of individual states) depend on the specific
transducer.  A function may have the finite decomposition property even when
some transducers realizing it fail the lemma's conditions.


\section{The Committed Prefix}
\label{sec:committed}

\begin{definition}[Committed prefix]
\label{def:committed}
For a function $\transFn \colon \srcStrings \to \tgtStrings$, the \defn{committed prefix} of an input $\srcStr \in \srcStrings$ is
\begin{equation}
  \committed(\srcStr) \defeq \lcp\!\set{\transFn(\srcStrp) : \srcStrp \succeq \srcStr}
\end{equation}
where $\lcp$ denotes the longest common prefix.
\end{definition}

Intuitively, $\committed(\srcStr)$ is the longest output prefix that is
\emph{guaranteed} after reading $\srcStr$, no matter what future input
arrives.

\begin{proposition}[Monotonicity]
\label{prop:committed-monotone}
$\committed$ is monotonically non-decreasing under the prefix order:
\[
  \srcStr \preceq \srcStrp
  \quad\Longrightarrow\quad
  \committed(\srcStr) \preceq \committed(\srcStrp).
\]
\end{proposition}
\begin{proof}
If $\srcStr \preceq \srcStrp$ then
$\{\transFn(\srcStr'') : \srcStr'' \succeq \srcStrp\} \subseteq
\{\transFn(\srcStr'') : \srcStr'' \succeq \srcStr\}$, so the lcp of the
smaller set is at least as long.
\end{proof}

\begin{proposition}[Committed prefix bounded by output]
\label{prop:committed-bound}
For every $\srcStr \in \srcStrings$:
$\committed(\srcStr) \preceq \transFn(\srcStr)$.
\end{proposition}
\begin{proof}
$\transFn(\srcStr)$ is an element of $\{\transFn(\srcStrp) : \srcStrp \succeq \srcStr\}$
(taking $\srcStrp = \srcStr$), so the lcp is at most $\transFn(\srcStr)$.
\end{proof}


\section{Recharacterizing the Decomposition}
\label{sec:recharacterize}

The committed prefix gives a clean function-level reading of the maximal
quotient and remainder.

\begin{proposition}[Maximal quotient via committed prefix]
\label{prop:maxQ-committed}
\begin{equation}
  \maxQuotient(\tgtStr) = \{ \srcStr \in \srcStrings : \tgtStr \preceq \committed(\srcStr) \}.
\end{equation}
\end{proposition}
\begin{proof}
By definition, $\srcStr \in \maxQuotient(\tgtStr)$ iff $\basisSet{\srcStr}
\subseteq \preCover(\tgtStr)$, i.e., for all $\srcStrp \succeq \srcStr$,
$\tgtStr \preceq \transFn(\srcStrp)$.  This holds iff $\tgtStr$ is a prefix
of every element of $\{\transFn(\srcStrp) : \srcStrp \succeq \srcStr\}$,
which is equivalent to $\tgtStr \preceq \lcp\{\transFn(\srcStrp) : \srcStrp
\succeq \srcStr\} = \committed(\srcStr)$.
\end{proof}

This yields a natural two-part decomposition of $\preCover(\tgtStr)$:
%
\begin{equation}
\preCover(\tgtStr) = \underbrace{\maxQuotient(\tgtStr)}_{\text{committed}}
  \;\sqcup\; \underbrace{\Remainder(\tgtStr)}_{\text{tentative}},
\end{equation}
%
where:
\begin{itemize}
\item \textbf{Committed set} $\maxQuotient(\tgtStr) = \{\srcStr : \tgtStr \preceq
  \committed(\srcStr)\}$: inputs after which the output is \emph{guaranteed} to
  start with $\tgtStr$, regardless of future input.

\item \textbf{Tentative set} $\Remainder(\tgtStr) = \{\srcStr : \tgtStr \preceq
  \transFn(\srcStr) \text{ but } \tgtStr \not\preceq \committed(\srcStr)\}$:
  inputs where the output \emph{happens to} start with $\tgtStr$, but some
  extension could steer the output away.  These are ``fragile'' coverers of
  $\tgtStr$.
\end{itemize}

\begin{remark}
The tentative set has a crisp reading as a \emph{monotonicity failure}: $\srcStr
\in \Remainder(\tgtStr)$ iff $\tgtStr \preceq \transFn(\srcStr)$ yet there
exists an extension $\srcStrp \succeq \srcStr$ with $\tgtStr \not\preceq
\transFn(\srcStrp)$.  That is, extending the input \emph{lost} the output
prefix $\tgtStr$---a violation of prefix monotonicity at $\srcStr$.
\end{remark}


\section{Two Independent Axes of Generalization}
\label{sec:axes}

The finite decomposition property requires both $\Quotient(\tgtStr)$ and
$\Remainder(\tgtStr)$ to be finite for all $\tgtStr$.  These are
\emph{independent} conditions on the function, each generalizing strict-prefix
monotonicity in a different direction.

\subsection{Axis A: Finite commitment boundary (finite quotient)}

Recall $\Quotient(\tgtStr) = \mathrm{pr}(\maxQuotient(\tgtStr))$ is the set of
minimal elements (in the prefix order) of the committed set.  These are the
\emph{earliest} input prefixes at which the function becomes irrevocably
committed to $\tgtStr$.  We call $\Quotient(\tgtStr)$ the \defn{commitment
frontier} for $\tgtStr$.

For the quotient to be finite, the committed set must have finitely many
minimal elements.  This can fail when the function allows \emph{arbitrarily
long delays} before commitment.

\begin{example}[Infinite quotient from unbounded delay]
\label{ex:infinite-Q}
Let $\srcAlphabet = \{a, b\}$, $\tgtAlphabet = \{c\}$, and define
$\transFn(\srcStr) = c^{\#_b(\srcStr)}$ (output one $c$ per $b$ in the
input).  Then:
\begin{itemize}
\item $\preCover(c) = \{\srcStr : \#_b(\srcStr) \geq 1\}$ = all strings containing at least one $b$.
\item $\committed(a^n) = \varepsilon$ for all $n$ (extending with only $a$'s produces no output).
\item $\committed(a^n b) = c$ (once a $b$ is read, at least one $c$ is guaranteed).
\item $\maxQuotient(c) = \{\srcStr : \srcStr \text{ contains at least one } b\}$.
\item $\Quotient(c) = \{a^n b : n \geq 0\} = \{b, ab, aab, aaab, \ldots\}$ --- \textbf{infinite}.
\end{itemize}
The function can ``stall'' for arbitrarily many input symbols (reading $a$'s)
without advancing the committed output.  There is no bound on how long
commitment to $c$ can be delayed.
\end{example}

\begin{definition}[Bounded delay]
\label{def:bounded-delay}
A function $\transFn \colon \srcStrings \to \tgtStrings$ has \defn{bounded delay}
if there exists a constant $k \geq 1$ such that for every $\srcStr \in
\srcStrings$ and every extension $\srcStrp \succeq \srcStr$ with $|\srcStrp| =
|\srcStr| + k$:
\begin{equation}
  |\committed(\srcStrp)| > |\committed(\srcStr)|
  \qquad\text{or}\qquad
  \committed(\srcStrp) = \committed(\srcStr) = \transFn(\srcStrp).
\end{equation}
That is, within every $k$ input symbols, the committed prefix either
\emph{advances} or the input has reached a string in the domain where the full
output is already committed.
\end{definition}

\begin{proposition}[Bounded delay implies finite quotient]
\label{prop:delay-quotient}
If $\transFn$ has bounded delay with constant $k$, then for every $\tgtStr \in \tgtStrings$:
\[
  |\Quotient(\tgtStr)| \leq |\srcAlphabet|^{k \cdot |\tgtStr|}.
\]
\end{proposition}
\begin{proof}[Proof sketch]
Every element of $\Quotient(\tgtStr)$ is a minimal string $\srcStr$ with
$\committed(\srcStr) \succeq \tgtStr$ but $\committed(\srcStrp) \not\succeq
\tgtStr$ for every proper prefix $\srcStrp \prec \srcStr$.  Since
$\committed(\varepsilon) = \lcp\{\transFn(\srcStrp) : \srcStrp \in
\srcStrings\}$ has some fixed length $\ell_0$, and $\committed$ advances by at
least one output symbol every $k$ input symbols, we need at most $k \cdot
(|\tgtStr| - \ell_0)$ input symbols for $\committed$ to reach length $|\tgtStr|$.
Hence every element of $\Quotient(\tgtStr)$ has length at most $k \cdot
|\tgtStr|$, and there are at most $|\srcAlphabet|^{k \cdot |\tgtStr|}$ such
strings.
\end{proof}

\begin{remark}[Connection to strict-prefix monotonicity]
Strict-prefix monotonicity is the special case $k = 1$.  If $\srcStr \prec
\srcStrp$ implies $\transFn(\srcStr) \prec \transFn(\srcStrp)$, then each
input symbol strictly advances the output, so
$|\committed(\srcStr \cdot \srcSym)| \geq |\committed(\srcStr)| + 1$ for
every $\srcSym \in \srcAlphabet$.  The bound recovers
$|\Quotient(\tgtStr)| \leq (|\srcAlphabet| + 1)^{|\tgtStr|}$ from the
paper's Proposition~B.2 (with the $+1$ accounting for elements
already in $\Quotient(\tgtStr)$ that persist to $\Quotient(\tgtStr\tgtSym)$).
\end{remark}

\begin{remark}[Connection to transducer theory]
Bounded delay is the function-level analog of the classical notion of
\emph{finite delay} in transducer theory~(Choffrut, 1977), which
characterizes when a rational function can be realized by a subsequential
transducer.  At the transducer level, it corresponds to the absence of
$\varepsilon$-output cycles---condition~(i) of the current Lemma~6.1.
\end{remark}


\subsection{Axis B: Finite tentative set (finite remainder)}

The remainder $\Remainder(\tgtStr)$ consists of inputs $\srcStr$ where
$\tgtStr \preceq \transFn(\srcStr)$ but $\tgtStr \not\preceq
\committed(\srcStr)$---the function's output happens to start with $\tgtStr$,
but future input could change this.

\begin{example}[Infinite remainder from unbounded non-monotonicity]
\label{ex:infinite-R}
Let $\srcAlphabet = \{a, b\}$, $\tgtAlphabet = \{c, d\}$, and define:
\[
  \transFn(\srcStr) = \begin{cases}
    c & \text{if } \srcStr = a^n \text{ for some } n \geq 1, \\
    d^{|\srcStr|} & \text{otherwise.}
  \end{cases}
\]
Then $\preCover(c) = \{a^n : n \geq 1\}$.  But $\committed(a^n) = \varepsilon$
for all $n \geq 1$, since $a^n b$ is an extension of $a^n$ with
$\transFn(a^n b) = d^{n+1}$ which does not start with $c$.  So
$\maxQuotient(c) = \emptyset$ (no input is committed to $c$), and
$\Remainder(c) = \{a, aa, aaa, \ldots\}$---\textbf{infinite}.

Every $a^n$ covers $c$ ``accidentally'': the output happens to be $c$, but
extending by $b$ steers it away.
\end{example}

\begin{definition}[Finite non-monotonicity]
\label{def:finite-nonmonotone}
A function $\transFn \colon \srcStrings \to \tgtStrings$ has \defn{finite
non-monotonicity} if for every $\tgtStr \in \tgtStrings$, the tentative set
\[
  \Remainder(\tgtStr) = \{\srcStr \in \srcStrings : \tgtStr \preceq \transFn(\srcStr) \text{ and } \tgtStr \not\preceq \committed(\srcStr)\}
\]
is finite.
\end{definition}

\begin{remark}[Connection to prefix monotonicity]
A prefix-monotone function has $\Remainder(\tgtStr) = \emptyset$ for all
$\tgtStr$ (paper, Proposition~5.1), since $\srcStr \preceq \srcStrp$ implies
$\transFn(\srcStr) \preceq \transFn(\srcStrp)$, so coverage is never
fragile.  Indeed, for prefix-monotone $\transFn$, we have
$\committed(\srcStr) = \transFn(\srcStr)$ for all $\srcStr$
(see \cref{prop:prefix-monotone-committed} below).

Finite non-monotonicity strictly generalizes this: the function
may fail prefix monotonicity at finitely many inputs per target prefix.
\end{remark}

\begin{proposition}
\label{prop:prefix-monotone-committed}
If $\transFn$ is prefix monotone, then $\committed(\srcStr) = \transFn(\srcStr)$ for all $\srcStr$.
\end{proposition}
\begin{proof}
By \cref{prop:committed-bound}, $\committed(\srcStr) \preceq \transFn(\srcStr)$.
For the reverse: prefix monotonicity gives $\transFn(\srcStr) \preceq
\transFn(\srcStrp)$ for all $\srcStrp \succeq \srcStr$.  So every element of
$\{\transFn(\srcStrp) : \srcStrp \succeq \srcStr\}$ extends $\transFn(\srcStr)$,
and their lcp satisfies $\committed(\srcStr) \succeq \transFn(\srcStr)$.
\end{proof}


\section{The Function-Level Finite Decomposition Lemma}
\label{sec:lemma}

\begin{lemma}[Function-level finite decomposition]
\label{lemma:function-level}
Let $\transFn \colon \srcStrings \to \tgtStrings$ be a function.  If $\transFn$
has bounded delay (\cref{def:bounded-delay}) and finite non-monotonicity
(\cref{def:finite-nonmonotone}), then $\transFn$ has the finite decomposition
property: for every $\tgtStr \in \tgtStrings$, both $\Quotient(\tgtStr)$ and
$\Remainder(\tgtStr)$ are finite.
\end{lemma}
\begin{proof}
Bounded delay implies $|\Quotient(\tgtStr)| \leq |\srcAlphabet|^{k \cdot |\tgtStr|}$ by \cref{prop:delay-quotient}.  Finite non-monotonicity directly asserts that $\Remainder(\tgtStr)$ is finite.
\end{proof}

\begin{remark}
The two conditions are independent:
\begin{itemize}
\item \Cref{ex:infinite-Q} has finite non-monotonicity ($\Remainder(c) = \emptyset$) but \emph{not} bounded delay, yielding an infinite quotient.
\item \Cref{ex:infinite-R} can be modified to have bounded delay but infinite non-monotonicity, yielding an infinite remainder.
\end{itemize}
\end{remark}



\section{Practical Implications}
\label{sec:practical}

\subsection{BPE tokenizers}

BPE detokenization is prefix-monotone (each token emits $\geq 1$ byte, and
extending the token sequence extends the byte output).  In fact it is
strict-prefix monotone when no token maps to $\varepsilon$.  So BPE satisfies
both conditions: bounded delay with $k = 1$, and $\Remainder = \emptyset$.

\subsection{PTB tokenizers}

The PTB tokenizer has context-dependent rules (e.g., ``\,\texttt{`\,`}\,''
$\to$ ``\,\texttt{"}\,'' at sentence boundaries, space insertion rules that
depend on surrounding tokens).  This can cause unbounded delay: arbitrarily
many input tokens can be consumed before the committed output advances, because
the output depends on future context.  This is the function-level reason why
PTB quotients are infinite.

\subsection{Connection to BFS termination}

The paper's algorithms decompose $\preCover(\tgtStr)$ via BFS over source
strings.  The BFS halts when every frontier element either:
\begin{enumerate}
\item \emph{commits} --- enters the committed set $\maxQuotient(\tgtStr)$
  (becomes a quotient element), or
\item \emph{exits} --- leaves $\preCover(\tgtStr)$ (not in the precover), or
\item \emph{resolves} --- is identified as a remainder element.
\end{enumerate}
Bounded delay ensures that the frontier cannot grow without bound (commitment
happens within bounded additional exploration).  Finite non-monotonicity ensures
that only finitely many frontier elements take path~(3).  Together, they
guarantee BFS termination---without reference to any particular transducer.


\section{The Committed Prefix as a Sequential Function}
\label{sec:sequential}

There is an additional structural observation about $\committed$.

\begin{proposition}
\label{prop:committed-sequential}
For any function $\transFn \colon \srcStrings \to \tgtStrings$, the committed
prefix $\committed \colon \srcStrings \to \tgtStrings$ is itself a
prefix-monotone function.
\end{proposition}
\begin{proof}
By \cref{prop:committed-monotone}, $\srcStr \preceq \srcStrp$ implies
$\committed(\srcStr) \preceq \committed(\srcStrp)$.
\end{proof}

So $\committed$ is the ``prefix-monotone part'' of $\transFn$---the largest
output that can be committed to incrementally.  The gap between
$\committed(\srcStr)$ and $\transFn(\srcStr)$ is what the remainder captures:
it consists of strings where the actual output overshoots the committed prefix.

This also gives an alternative characterization of the decomposition:
\begin{equation}
  \preCover(\tgtStr) = \underbrace{\preCover_{\committed}(\tgtStr)}_{\text{committed part}}
  \;\cup\;
  \underbrace{\{\srcStr : \tgtStr \preceq \transFn(\srcStr),\;
    \tgtStr \not\preceq \committed(\srcStr)\}}_{\text{remainder}},
\end{equation}
where $\preCover_{\committed}(\tgtStr) = \{\srcStr : \tgtStr \preceq
\committed(\srcStr)\}$ is the precover of $\tgtStr$ under the prefix-monotone
function $\committed$.  Since $\committed$ is prefix-monotone, this precover
is automatically a cylinder, and its quotient is exactly $\Quotient(\tgtStr)$.


\section{Necessity of Bounded Delay}
\label{sec:necessity}

\Cref{sec:axes} established that bounded delay is \emph{sufficient} for finite
quotients.  Is it \emph{necessary}?  The answer depends on whether we consider
arbitrary functions or restrict to FST-realizable ones.

\subsection{General functions: bounded delay is not necessary}

\begin{proposition}
\label{prop:not-necessary-general}
There exists a prefix-monotone function $\transFn \colon \srcStrings \to
\tgtStrings$ such that $\Quotient(\tgtStr)$ is finite for every $\tgtStr$ but
$\transFn$ does not have bounded delay.
\end{proposition}
\begin{proof}
Let $\srcAlphabet = \{a, b\}$, $\tgtAlphabet = \{c\}$, and define
$\transFn(\srcStr) = c^{\lfloor \log_2(|\srcStr| + 1) \rfloor}$.  That is,
the output length is the floor of $\log_2$ of one more than the input length.
This function depends only on $|\srcStr|$, and since $\lfloor \log_2(n+1)
\rfloor \leq \lfloor \log_2(m+1) \rfloor$ whenever $n \leq m$, the function
is prefix monotone (longer inputs produce longer-or-equal outputs).  By
\cref{prop:prefix-monotone-committed}, $\Remainder(\tgtStr) = \emptyset$ for
all $\tgtStr$.

\medskip\noindent\textbf{Committed prefix.}
$\committed(\srcStr) = c^{\lfloor \log_2(|\srcStr| + 1) \rfloor}$, since the
minimum output length over all extensions of $\srcStr$ is achieved by $\srcStr$
itself (shorter extensions have shorter output).

\medskip\noindent\textbf{Finite quotients.}
$\maxQuotient(c^n) = \{ \srcStr : |\srcStr| \geq 2^n - 1 \}$, so
$\Quotient(c^n) = \srcAlphabet^{2^n - 1}$ (all strings of length exactly
$2^n - 1$).  This is finite: $|\Quotient(c^n)| = 2^{2^n - 1}$.

\medskip\noindent\textbf{No bounded delay.}
Between input lengths $2^n - 1$ and $2^{n+1} - 2$, the committed output stays
at $c^n$.  This stalling interval has length $2^n - 1$, which grows without
bound.  So no constant $k$ satisfies \cref{def:bounded-delay}.
\end{proof}

\begin{remark}
This counterexample is not FST-realizable.  The output schedule requires
counting to $2^n$ between successive output symbols, which exceeds the
capacity of any finite-state device.
\end{remark}

\subsection{FST-realizable functions: bounded delay is necessary}

For functions realized by finite-state transducers, the situation is different.
The finite state space forces a dichotomy: either the committed prefix advances
within a bounded number of steps, or it can stall indefinitely---and
indefinite stalling creates infinite quotients.

\begin{proposition}
\label{prop:necessary-fst}
Let $\transFn \colon \srcStrings \to \tgtStrings$ be a function realized by an
input-deterministic FST $\transST$ with $N$ states.  If $\transFn$ does not
have bounded delay with constant $k = N$, then there exists
$\tgtStr \in \tgtStrings$ such that $\Quotient(\tgtStr)$ is infinite.
\end{proposition}
\begin{proof}
Suppose bounded delay fails with $k = N$.  Then there exist $\srcStr \in
\srcStrings$ and symbols $a_1, \ldots, a_N \in \srcAlphabet$ such that
$\committed(\srcStr \cdot a_1 \cdots a_i) = \committed(\srcStr)$ for all
$i = 0, \ldots, N$, and $\committed(\srcStr \cdot a_1 \cdots a_N) \neq
\transFn(\srcStr \cdot a_1 \cdots a_N)$.

Since $\transST$ is input-deterministic, processing $\srcStr$ then
$a_1, \ldots, a_N$ visits $N + 1$ states.  By pigeonhole, some state $\state$
is visited twice: $\state$ is reached after $\srcStr \cdot a_1 \cdots a_j$
and again after $\srcStr \cdot a_1 \cdots a_{j'}$ for some $0 \leq j < j'
\leq N$.

\medskip\noindent\textbf{Claim:} The cycle word $w = a_{j+1} \cdots a_{j'}$
produces $\varepsilon$ output.

Since the transducer is input-deterministic, the output produced along
$a_{j+1}, \ldots, a_{j'}$ from $\state$ is the same string $\tgtStrp$ every
time the cycle is traversed.  Suppose $\tgtStrp \neq \varepsilon$.  Then
traversing the cycle $m$ times from $\state$ produces output $\tgtStrp^m$.
Consider the string $\srcStr' = \srcStr \cdot a_1 \cdots a_j$ (which reaches
$\state$).  The output of $\transFn(\srcStr' \cdot w^m \cdot z)$ for any
continuation $z$ begins with $\transFn(\srcStr') {|}_{\text{prefix}} \cdot
\tgtStrp^m$ (the output accumulated before the cycle, plus $m$ copies of
$\tgtStrp$, plus whatever $z$ produces).  Therefore
$|\committed(\srcStr' \cdot w^m)| \geq |\committed(\srcStr')| + m \cdot
|\tgtStrp|$, which grows with $m$.  But $\committed(\srcStr' \cdot w) =
\committed(\srcStr \cdot a_1 \cdots a_{j'}) = \committed(\srcStr)$
(by assumption that committed stalls).  This contradicts $|\tgtStrp| \geq 1$.
Hence $\tgtStrp = \varepsilon$.

\medskip\noindent\textbf{The $\varepsilon$-output cycle creates an infinite quotient.}

Let $\srcStr' = \srcStr \cdot a_1 \cdots a_j$, so $\srcStr'$ reaches $\state$
with accumulated output $\tgtStr_0$.  Since $w$ produces $\varepsilon$ and
returns to $\state$, we have $\transFn(\srcStr' \cdot w^m \cdot z) =
\transFn(\srcStr' \cdot z)$ for all $m \geq 0$ and $z \in \srcStrings$.  In
particular, $\committed(\srcStr' \cdot w^m) = \committed(\srcStr')$ for all
$m$.

Since $\committed(\srcStr') \neq \transFn(\srcStr')$ (from the bounded-delay
failure assumption), there exists a symbol $\tgtSym$ such that
$\tgtStr_1 \defeq \committed(\srcStr') \cdot \tgtSym$ satisfies
$\committed(\srcStr') \prec \tgtStr_1 \preceq \transFn(\srcStr' \cdot z_0)$
for some continuation $z_0$.  Now consider the elements of $\Quotient(\tgtStr_1)$.
Any $\srcStr''$ with $\committed(\srcStr'') \succeq \tgtStr_1$ must extend
$\srcStr'$ past the $\varepsilon$-cycle (i.e., it must eventually take a
non-cycle transition from $\state$).  Concretely, for each $m \geq 0$, the
string $\srcStr' \cdot w^m \cdot b$ (where $b$ is the first symbol of $z_0$)
reaches the same post-cycle state with $\committed \succeq \tgtStr_1$.  These
strings are pairwise prefix-incomparable---$\srcStr' \cdot w^m \cdot b$ and
$\srcStr' \cdot w^{m'} \cdot b$ differ at position $|\srcStr'| + (m \wedge m')
\cdot |w| + 1$ (one has $b$, the other has $w$'s first symbol, and $b$ differs
from $w$'s first symbol because $b$ exits the cycle while $w$'s first symbol
stays on it).  Hence $\Quotient(\tgtStr_1)$ contains infinitely many minimal
elements.
\end{proof}

\begin{corollary}
\label{cor:fst-equivalence}
For FST-realizable functions, bounded delay is both necessary and sufficient
for finite quotients:
\[
  \text{$\Quotient(\tgtStr)$ finite for all $\tgtStr$}
  \quad\Longleftrightarrow\quad
  \text{$\transFn$ has bounded delay (with $k \leq N$)}.
\]
\end{corollary}

\begin{remark}
The $\varepsilon$-output cycle in the proof is a structural feature of the
function, not an artifact of the chosen transducer.  If $\transFn$ does not
have bounded delay, then \emph{every} input-deterministic FST realizing
$\transFn$ must contain such a cycle (by the pigeonhole argument).
\end{remark}

\begin{remark}[Bounded vs.\ finite delay]
The counterexample in \cref{prop:not-necessary-general} has \emph{finite
delay}---$\committed$ always advances in finite time, just not within a uniform
bound---but not \emph{bounded delay}.  For FST-realizable functions, the
distinction collapses: by \cref{prop:necessary-fst}, finite delay implies
bounded delay (with $k \leq N$), because an FST with finite delay cannot have
$\varepsilon$-output cycles.
\end{remark}


\section{Summary}
\label{sec:summary}

\begin{center}
\begin{tabular}{@{}lcc@{}}
\toprule
\textbf{Condition on $\transFn$} & \textbf{Finite $\Quotient$?} & \textbf{Finite $\Remainder$?} \\
\midrule
Strict-prefix monotone & Yes ($\leq (|\srcAlphabet|+1)^{|\tgtStr|}$) & Yes ($= \emptyset$) \\
Prefix monotone & Not guaranteed & Yes ($= \emptyset$) \\
Bounded delay & Yes ($\leq |\srcAlphabet|^{k|\tgtStr|}$) & Not guaranteed \\
Finite non-monotonicity & Not guaranteed & Yes (by definition) \\
\midrule
Bounded delay $+$ finite non-mono. & Yes & Yes \\
\bottomrule
\end{tabular}
\end{center}

\medskip\noindent
For FST-realizable functions, bounded delay is an \textbf{equivalence}:
$\Quotient(\tgtStr)$ is finite for all $\tgtStr$ if and only if $\transFn$ has
bounded delay (\cref{cor:fst-equivalence}).  For general functions, it is only
sufficient (\cref{prop:not-necessary-general}).


\section{Remaining Open Questions}
\label{sec:open}

\begin{enumerate}
\item \textbf{Necessity of finite non-monotonicity.}
  We showed bounded delay is necessary (for FSTs) for finite quotients.
  Is finite non-monotonicity necessary for finite remainders?  That is: if
  $\Remainder(\tgtStr)$ is finite for all $\tgtStr$, must the tentative set
  be finite?  (This is tautological as stated; a non-trivial version would
  seek a characterization in terms of more elementary properties of $\transFn$.)

\item \textbf{Tighter bounds.}
  The bound $|\srcAlphabet|^{k|\tgtStr|}$ is crude.  Can it be sharpened, e.g.,
  using the structure of the committed prefix's ``increments''?

\item \textbf{Decidability.}
  Given an FST realizing $\transFn$, is bounded delay decidable?  For rational
  functions, finite delay is known to be decidable (Choffrut 1977), and
  our \cref{prop:necessary-fst} shows that finite delay = bounded delay for
  FSTs, so the answer is likely yes.  Is finite non-monotonicity decidable?
\end{enumerate}


\end{document}
